\documentclass[10pt,a4paper]{amsart}
\usepackage[margin=19mm]{geometry}
\usepackage{amsmath,amssymb,mathtools}
\usepackage[scr=rsfs,cal=euler]{mathalfa}
\usepackage{xfrac}
\usepackage[unicode,hidelinks,bookmarks=false]{hyperref}
\newcommand{\ie}{i.e.\ }
\newcommand{\cf}{cf.\ }
\newcommand{\resp}{respectively\ }
\newcommand{\Subsection}{\subsection}
\providecommand{\nobreakdash}{\nobreak\hbox{-}}
\setlength{\emergencystretch}{2em}
\multlinegap0pt
\usepackage{amssymb}
\usepackage{enumerate}
\newtheorem{proposition}{Proposition}
\newcommand{\Ker}{\mathrm{Ker}}
\newcommand{\Lie}{\mathrm{Lie}}
\title{A Buium--Coleman bound for the Zilber--Pink conjecture for curves inside abelian varieties}
\author{Netan Dogra, Sudip Pandit}
\date{Source: arXiv:2608.24474v1 (CC BY 4.0)}
\begin{document}
\maketitle

\noindent\emph{Source and licence.} A Buium--Coleman bound for the Zilber--Pink conjecture for curves inside abelian varieties. Version arXiv:2608.24474v1 (CC BY 4.0), distributed by arXiv under the Creative Commons Attribution 4.0 International licence, http://creativecommons.org/licenses/by/4.0/, as recorded in the arXiv licence metadata for this exact version and shown on the abstract page https://arxiv.org/abs/2608.24474v1. Full licence notice: \texttt{licenses/arxiv-2608.24474-CC-BY-4.0.txt}.\par\smallskip

\noindent\textit{Editorial scope.} Counted unit: the article's proposition prop:finite\_2 (source0.tex lines 515--517). The article's complete original proof of it is delivered with it (source0.tex lines 518--557, one \texttt{proof} environment of 40 lines). No mathematics is written, repaired or re-derived: the statement and the proof are the article's own lines, in the article's own notation, and no retained line is substituted or re-flowed.  No further source-local context is needed: every cross-reference of every retained line resolves inside the two delivered spans.  Nothing was omitted from any retained span, so the package carries no omission record.  No retained line contains a citation, so the wrapper carries no bibliography and no bibliography entry.  No retained line contains a figure environment, an image-inclusion command or drawing code, so no image is delivered and the package declares no asset dependency.  Editorial formatting only, in addition: an AMS \texttt{amsart} preamble that restates the article's own declarations and macros used by the retained lines, namely the environments \texttt{proposition} on separate counters, together with the standard packages those lines need.  The retained environments print in this document as Proposition 1 in order of appearance, whereas the article prints the same items with the numbering of its own full document.  The complete native source of the article as distributed in the arXiv e-print for this exact version is bundled unchanged as \texttt{sources/208480/source0.tex}.
\begin{proposition}\label{prop:finite_2}
For $A$ and $W$ as above, the number of elements of $A(R/p^2 R)/pA(R/p^2 R)+B(R/p^2 R)$ which map to torsion points in $A/B$, for $B<A$ with $(A/B)[p]\simeq W$, is at most $p^{g_A}$.
\end{proposition}

\begin{proof}
Suppose $B_1 $ and $B_2 $ are abelian subvarieties with $(A/B_i )[p]\simeq W$, that $P_i \in A(R)$ is such that $P_i$ maps to a torsion point in $A/B_i$, and that the classes of $P_i$ in $H^1 (K,W)$ are the same.

Let $f:A'\to A$ be the finite \'etale cover corresponding to $W$. Then we have an injection
\[
A(R)/f(A'(R))+pA(R)\hookrightarrow H^1 (K,W).
\]
On the other hand the map
\[
A(R)/pA(R)+B_1 (R)+B_2 (R) \to H^1 (K,W)
\]
is the composite of the inclusion
\[
A(R)/pA(R)+B_1 (R)+B_2 (R) \hookrightarrow A(R)/A'(R)
\]
induced by the map
\[
B_1 \times B_2 \to A'
\]
(such a map exists because under the map $B_1 \times B_2 \to A$, the image of $(B_1 \times B_2 )[p]$ is $\Ker (A[p]\to W)$) with the injection
\[
A(R)/A'(R)\to H^1 (K,W)
\]
coming from the exact sequence
\[
0\to M\to A' (\overline{K})\to A(\overline{K})\to 0.
\]
We deduce that the images of $P_1$ and $P_2 $ in $A(R)/pA(R)+B_1 (R)+B_2 (R)$ are the same. Hence the images in $A(R/p^2 R)/pA(R/p^2 R)+B_1 (R/p^2 R)+B_2 (R/p^2 R)$ are the same. The result follows from the isomorphisms
\[
pA(R/p^2 R)+B_1 (R/p^2 R)+B_2 (R/p^2 R)\simeq pA(R/p^2 R)+B_i (R/p^2 R)
\]
for $i=1,2$. Indeed, since $B_1 [p]=B_2 [p]$ as subspaces of $A[p]$, we deduce $\Lie (B_{1,k})=\Lie (B_{2,k})$ as subspaces of $\Lie (A_k )$, since the morphisms of group schemes
\[
B_i [p]_k \hookrightarrow B_{i,k}
\]
induce isomorphism on Lie algebras. Hence the images of $J^1 B_{1,k}$ and $J^1 B_{2,k}$ in $J^1 A_k (k)/pJ^1 A_k (k)$ are the same.
%\[
%pA(R/p^2 R)+B_1 (R/p^2 R)=pA(R/p^2 R)+B_2 (R/p^2 R).
%\]
\end{proof}

\end{document}
